\documentclass[11pt]{article}

\usepackage[T1]{fontenc}
\usepackage{lmodern}
\usepackage{xcolor}
\pagecolor{white}

\usepackage[margin=1in]{geometry}
\usepackage{amsmath,amssymb,amsthm}
\usepackage{booktabs}
\usepackage{microtype}

\newtheorem{law}{Law}
\newtheorem{remark}{Remark}

\newcommand{\bthree}{b_3}
\newcommand{\flux}{F}
\newcommand{\hh}{H}
\newcommand{\tauu}{\tau}

\title{Mitchell's Laws of Duality Covariance: the covariance table of the flux landscape}
\author{Rob Mitchell}
\date{}

\begin{document}
\maketitle

\begin{center}\itshape
``Don't ask a complex system for its values --- ask for its constraints, its count, and its
deficit. Everything else is covariance, not knowledge.''
\end{center}

\begin{abstract}
The constraint-first theorem proves that the D3 tadpole $N=\flux^\top J\hh$ is the
\emph{only} exact duality-invariant of the type IIB flux, so the invariant ring is
$\mathbb{C}[N]$ and no exact invariant selector can single out a vacuum. That theorem
says \emph{what survives}; it does not say \emph{how the surviving observables move}
under the duality that fixes them. This paper measures that motion. On $100$ primitive
supersymmetric flux vacua of the KS truncation ($h^{1,2}=2$), with every measurement
gated by a control that reproduces a known value, we determine the transformation law
of the flux superpotential $W$ and the string coupling $g_s$ under the two generators
of $\mathrm{SL}(2,\mathbb{Z})$:
\begin{itemize}
  \item \textbf{Law I (scaling).} $W(d\flux,d\hh)=d\,W$, $N(d\flux,d\hh)=d^2N$; the
    content $d=\gcd(\flux,\hh)$ is a bookkeeping label, not a selector.
  \item \textbf{Law II (magnitude).} $|W(-\hh,\flux;-1/\tauu)|=|W(\flux,\hh;\tauu)|$,
    ratio $1.000000$ (log residual $0.0000$) on $100/100$ vacua.
  \item \textbf{Law IIb (phase).} $\arg(W_S/W)=-\arg\tauu$ exactly (fit slope
    $-1.0000$, residual $0.0000$), so $W_S = W\,\bar\tauu/|\tauu|$: S-duality acts on
    the gauge-invariant superpotential by a pure phase.
  \item \textbf{Law III ($g_s$).} $g_s=1/\mathrm{Im}\,\tauu$ transforms as
    $g_s\to|\tauu|^2 g_s$ under $S$ and $g_s\to g_s$ under $T$ (ratios $1.000000$).
  \item \textbf{Well-definedness.} the $S$-dual of a stable vacuum is stable
    ($|DW|=1.8\times10^{-12}$ at the transformed point, control reproduced the stored
    $3.2\times10^{-12}$).
\end{itemize}
The one place a law \emph{fails} to close --- $|W|$ under $T$ --- is reported as a
finding, not suppressed. Together these rows are the \emph{covariance table}: the
operational statement of ``one theory described twice,'' and the empirical half of the
unification programme. Every law carries a named falsifier; the measurement code and
data are public in \texttt{github.com/robmitchellengineering-hub/pacland}.
\end{abstract}

\section{Introduction}

The constraint-first programme's central theorem \cite{constraint-first} identifies the
exact content of the flux sector: the invariant ring of the flux pair under
$\mathrm{Sp}(\bthree,\mathbb{Z})\times\mathrm{SL}(2,\mathbb{Z})$ is generated by the
tadpole alone, so the exact structure of a theory of everything is \emph{constraints
plus labels}, nothing else. The proof is symbolic and machine-checked.

That theorem fixes the \emph{invariants}. It leaves open a different, equally sharp
question: how do the \emph{non-invariant} observables --- the superpotential value
$W$, the string coupling $g_s$, the moduli --- transform under a duality that leaves
the tadpole fixed? The answer is the \emph{covariance table}: a closed-form law
$O\mapsto f_g(O,\tau,z)$ for each observable $O$ and each duality generator $g$. A
duality is ``one theory described twice'' only if that table exists and closes; where
it fails to close, the failure is itself a finding.

This paper records the first complete rows of that table, measured on real
supersymmetric vacua with a control attached to every number.

\section{Setup and method}

Let $X$ be a Calabi--Yau threefold, $\flux,\hh\in H^3(X,\mathbb{Z})$ the RR and NS--NS
three-form fluxes, $J$ the symplectic form, and $\tauu=C_0+i e^{-\phi}$ the
axio-dilaton ($g_s=e^{\phi}=1/\mathrm{Im}\,\tauu$). The duality generators act as
\begin{equation}
  S:\ (\flux,\hh)\to(-\hh,\flux),\quad \tauu\to-1/\tauu;
  \qquad
  T:\ (\flux,\hh)\to(\flux,\flux+\hh),\quad \tauu\to\tauu+1,
\end{equation}
at fixed moduli (the degree-0 slice). Both preserve the tadpole by antisymmetry of
$J$.

\textbf{Method.} Each vacuum is a row of the committed corpus
(\texttt{data/h12\_*\_stats.json}); we take the $100$ primitive ($d=1$) supersymmetric
vacua of the KS truncation at $h^{1,2}=2$, $N_{\rm max}=0$. The superpotential is the
gauge-invariant value $W(\flux,\tauu)$ from \texttt{jaxvacua} (0.1.2). \textbf{Every
measurement is gated by a control that reproduces a known value:} the $|W|$ rows
reproduce the published Law II on the same code path, and the well-definedness row
reproduces the stored $|DW|$ of its own vacuum. A control that fails suppresses the
row, it does not retry. A job whose exit code a shell pipe could mask is never
reported as a pass --- the first G2 run was exactly such a false pass (a missing
$\tauu_c$ argument), and it is fixed and re-run here.

\section{The laws}

\begin{law}[scaling]\label{law:scaling}
For a superpotential linear in the flux,
\begin{equation}
  W(d\flux,d\hh)=d\,W,\qquad DW(d\flux,d\hh)=d\,DW,\qquad
  N(d\flux,d\hh)=d^2 N,\qquad \gcd(d\flux,d\hh)=d\gcd(\flux,\hh).
\end{equation}
So every SUSY vacuum at moduli $z$ has a SUSY vacuum at the \emph{same} $z$ with flux
scaled by $d$; the content $d$ changes $W\to dW$ and $N\to d^2N$ and nothing else. It
is a lattice label, not a selector.
\end{law}

\textit{Verification.} Symbolic, exact ($5/5$, \texttt{validate\_flux\_scaling\_lemma.py});
against the library on a real vacuum, $|W(2\flux)|/|W(\flux)|=2.0$ at tolerance
$10^{-10}$. \textit{Scope.} $W$ linear in flux (tree level). \textit{Falsifier.} a DW
zero set not preserved under flux scaling.

\begin{law}[S-duality invariance of $|W|$]\label{law:mag}
$|W(-\hh,\flux;-1/\tauu)|=|W(\flux,\hh;\tauu)|$, equivalently the realised law is
$|W_S|/|W|=|\tauu|^0$.
\end{law}

\textit{Verification.} $100/100$ primitive SUSY vacua, free fit of both exponents in
$|W_S|/|W|=e^{a}|\tauu|^{b}$ returns $a=b=0$ with log residual $0.0000$. \textit{Scope.}
$h^{1,2}=2$, $N_{\rm max}=0$, magnitude only. \textit{Falsifier.} a log residual
$>0.10$.

\begin{law}[phase]\label{law:phase}
$\arg(W_S/W)=-\arg\tauu$ exactly, i.e.\ $W_S=W\,\bar\tauu/|\tauu|$: the gauge-invariant
superpotential is S-covariant by a pure phase.
\end{law}

\textit{Verification.} same $100$ vacua; fit slope $-1.0000$, intercept $0.0000$,
residual $0.0000$; the magnitude of the same ratio reproduces Law~\ref{law:mag}
($1.000000$). \textit{Falsifier.} a phase that scatters away from $-\arg\tauu$.

\begin{law}[string-coupling covariance]\label{law:gs}
$g_s=1/\mathrm{Im}\,\tauu$ transforms as $g_s\to|\tauu|^2 g_s$ under $S$ and
$g_s\to g_s$ under $T$.
\end{law}

\textit{Verification.} ratios $g_s'/(|\tauu|^2 g_s)=1.000000$ and $g_s'/g_s=1.000000$
on $100/100$ vacua. \textit{Scope.} this is definitional --- $g_s$ is read off $\tauu$
--- so a deviation is a measurement bug, not a physics finding (it exposed an earlier
``$g_s\to 3.9\times10^7$'' error). \textit{Falsifier.} $g_s'\neq|\tauu|^2 g_s$ under
$S$.

\section{Well-definedness}

The laws above are necessary but not sufficient for a duality to be a \emph{symmetry
of the vacuum set}: the $S$-dual of a stable vacuum must itself be stable. On the
first primitive SUSY vacuum ($h^{1,2}=2$, $N_{\rm flux}=57$), the control reproduces
the stored $|DW|=3.19\times10^{-12}$ against $4.35\times10^{-12}$, and the transformed
point satisfies $|DW|(-\hh,\flux;-1/\tauu)=1.83\times10^{-12}\le10^{-6}$. \textbf{The
S-dual of a stable vacuum is stable}; C3's well-definedness holds at degree 0.

\section{Where a law fails to close}

Under $T$, $|W_T|/|W|$ does \emph{not} close to a clean law: mean $4.9\times10^6$,
standard deviation $4.9\times10^7$, fit residual $0.22$. The mean is carried by rows
where $|W|\approx0$ (the no-scale locus, where the ratio is ill-defined); the generic
rows show a rough $|W_T|/|W|\approx0.6\,|\tauu|$. This is recorded as a finding, not a
law: the gauge-invariant $|W|$ is S-covariant (Laws II, IIb) but not T-covariant in a
closed form. Whether a clean T-law exists away from $W=0$ is the named follow-up.

\section{The unification reading}

The theorem \cite{constraint-first} shows the orbit space is a \emph{point} per tadpole
(over $\mathbb{R}$, $\mathrm{Sp}(\bthree)$ is transitive on each level set), so
observables cannot be \emph{coordinates} on the orbit space. The covariance table is
what remains: the observables are covariant \emph{functions} on the orbit, and a
duality is a symmetry of the theory exactly when every row of the table closes. Laws
I--III and well-definedness close the S-generator rows; the T-row of $|W|$ is open.
This is the operational meaning of ``one theory described twice,'' and it is the
empirical half of the unification programme.

\section{Falsifiers}

Each law carries its own falsifier above. The programme-level falsifiers are unchanged
\cite{constraint-first}: (i) a moduli-independent, duality-invariant flux quantity not
a function of $\flux^\top J\hh$; (ii) an exact UV relation among the Standard-Model
parameters beyond the known GUT collapses; (iii) the invariant structure changing as
more moduli are kept (the truncation diagnostic, which returns \emph{interior}, i.e.\ a
physical hierarchy, not a truncation artefact).

\begin{thebibliography}{9}
\bibitem{constraint-first} R.~Mitchell, \emph{Mitchell's Completeness Theorem: no exact
invariant selector in the flux landscape}, companion paper, this series.
\end{thebibliography}

\end{document}
